\documentclass[journal]{IEEEtran}
\usepackage{amsmath,amssymb}
\usepackage{booktabs}
\usepackage{graphicx}
\usepackage{hyperref}
\usepackage{multirow}

\begin{document}

\title{Safety Without Understanding: Why Conservative Value Functions Collapse to Behavioral Cloning in Combinatorial Selection}
\author{}
\maketitle

\begin{abstract}
When learned value functions guide search in combinatorial selection, two capabilities determine policy quality: \textit{composition safety} (avoiding structurally broken selections) and \textit{context-awareness} (selecting elements that interact well given the specific situation). We show that existing methods trade one for the other. In Heroes of the Storm draft selection, three win probability models spanning 56.5--57.9\% test accuracy disagree on 90\%+ of decisions. Feature engineering with synthetic out-of-distribution augmentation reduces degenerate compositions from 76\% to 26.5\% while producing context-aware drafts with positive counter responsiveness (+0.06), strong synergy exploitation (+0.80), and 90 distinct heroes across 200 drafts. Conservative Q-Learning (CQL) appears to solve composition safety more completely (9\% degenerate), but rich evaluation reveals it collapses to behavioral cloning: negative counter/synergy deltas, half the hero diversity, and 62\% agreement with the behavior policy. The behavioral cloning baseline itself achieves 99.5\% healer selection and 1\% degenerate rate, confirming that composition safety was never the hard problem. MCTS with the augmented value function achieves 0.626 average win probability (vs.\ 0.602 without augmentation), combining deeper search with a well-calibrated value function. Composition safety and context-awareness are orthogonal capabilities: algorithmic pessimism provides the former by mimicking human behavior; domain-structured features with OOD augmentation provide both.
\end{abstract}

\section{Introduction}

When a learned value function guides a search or planning algorithm, a natural assumption is that higher held-out accuracy produces better policies. We show that this assumption fails in a specific and instructive way, and that the failure reveals two orthogonal axes of policy quality that existing methods cannot optimize simultaneously.

The problem arises in combinatorial team selection: choosing a set of elements where both individual quality and interactions between elements determine outcomes. Standard evaluation measures aggregate prediction accuracy on held-out data from the same distribution. But when the value function is used for search, the optimizer discovers states outside the training distribution, and the value function's behavior on those states determines whether the resulting policy makes reasonable selections.

We study this in Heroes of the Storm (HotS) draft selection, a sequential two-player game where teams alternately ban and pick characters (``heroes'') from a pool of 90. The resulting 5v5 team compositions determine roughly 55--58\% of game outcomes. The domain is well-suited because the state space requires generalization ($\binom{90}{5}^2 \approx 10^{13}$ possible matchups), training data concentrates on reasonable compositions, and the gap between individual element quality and compositional quality is large and measurable.

We decompose draft quality into two orthogonal axes:
\begin{itemize}
    \item \textbf{Composition safety:} Does the team include the right structural roles (healer, frontline, ranged damage)? Measured by healer rate, degenerate composition rate.
    \item \textbf{Context-awareness:} Does the team counter the opponent's picks and exploit internal synergies? Measured by pairwise counter responsiveness, synergy exploitation, hero diversity, and agreement with behavioral cloning.
\end{itemize}

Existing methods provide one axis or the other. Behavioral cloning of human drafts trivially solves composition safety (99.5\% healer rate, 1\% degenerate) because humans always draft balanced teams. Conservative Q-Learning~\cite{kumar2020conservative} achieves 9\% degenerate rate but collapses to behavioral cloning on rich evaluation metrics: negative counter/synergy deltas, half the hero diversity, 62\% agreement with the behavior policy. Value function search with enriched features produces context-aware drafts (positive counter/synergy, 90 distinct heroes, 4.4\% GD agreement) but requires feature engineering and synthetic OOD augmentation to also achieve composition safety.

Our contributions:
\begin{enumerate}
    \item We demonstrate that value functions with a 1.4pp accuracy difference disagree on 90\%+ of decisions and produce drastically different compositions.
    \item We show that feature engineering plus synthetic augmentation for unseen compositions produces both safe \textit{and} context-aware drafts (healer 94.5\%, degenerate 26.5\%, counter +0.06, synergy +0.80, 90 heroes).
    \item We provide the first empirical demonstration that CQL collapses to behavioral cloning in combinatorial selection, achieving safety by mimicking humans rather than understanding composition.
    \item We decompose draft quality into orthogonal axes (safety vs.\ context-awareness) with quantitative metrics that distinguish genuine reasoning from behavioral mimicry.
    \item We show MCTS with an augmented value function achieves the highest overall performance (0.626 avg WP, 92.5\% win rate), combining deeper search with a well-calibrated value function.
\end{enumerate}

\section{Background and Related Work}

\subsection{MOBA Draft Prediction}

Chen et al.~\cite{chen2018drafting} proposed MCTS-based draft recommendation for Dota~2, using a win rate predictor as the value function for tree search. Their system uses multi-hot hero features and evaluates drafts with the same model used for search, without examining how value function features affect policy quality or out-of-distribution behavior. Gourdeau and Archambault~\cite{gourdeau2021discriminative} trained discriminative neural networks for hero selection in professional HotS and Dota~2. We reproduce their architecture and show it exhibits OOD overestimation, scoring 5-tank compositions at 52.8\% win probability (rating them as favored). Summerville et al.~\cite{summerville2016draft} modeled professional HotS draft sequences as behavioral cloning with LSTMs. DraftRec~\cite{lee2022draftrec} proposes a transformer-based personalized recommendation system that explicitly restricts hero suggestions to combinations observed in the training distribution, implicitly addressing the OOD problem we characterize but at the cost of novel strategy discovery. Hanke and Chaimowicz~\cite{hanke2017recommending} applied association rules to hero recommendations, a purely statistical approach without learned value functions. Pobiedina et al.~\cite{pobiedina2013draft} studied team composition effects on match outcomes in Dota~2, establishing that composition matters beyond individual hero strength.

No prior work examines how value function features affect policy quality, measures the composition gap between accuracy and policy behavior, tests whether CQL provides genuine compositional reasoning vs.\ behavioral cloning, or decomposes draft quality into safety and context-awareness.

\subsection{Offline RL and OOD Overestimation}

Conservative Q-Learning (CQL)~\cite{kumar2020conservative} addresses OOD value overestimation by adding a penalty term that pushes down Q-values for actions underrepresented in the training data. BEAR~\cite{kumar2019stabilizing} constrains the learned policy to remain close to the behavior policy via distributional matching. Levine et al.~\cite{levine2020offline} provide a comprehensive survey of offline RL methods and the OOD overestimation problem. Lyu et al.~\cite{lyu2022mildly} propose Mildly Conservative Q-Learning (MCQ), addressing the tendency of CQL to be overly conservative in practice, a finding consistent with our observation that CQL collapses to behavioral cloning.

Our contribution to this literature is empirical: we provide the first demonstration that CQL's pessimism eliminates context-awareness in combinatorial selection, achieving composition safety by the same mechanism as behavioral cloning rather than by learning compositional reasoning.

\subsection{Value Function Design for Planning}

The sensitivity of planning algorithms to value function quality has been studied in game-playing~\cite{silver2016mastering} and robotic manipulation~\cite{kalashnikov2018qtopt}. Our work provides a controlled experiment where three value functions with near-identical accuracy produce policies that disagree on over 90\% of decisions, with a quantitative decomposition of the disagreement into orthogonal safety and context-awareness axes.

Our earlier work~\cite{segan2026composition} identified the composition gap and proposed feature engineering plus synthetic augmentation as mitigations. Here we show that algorithmic pessimism (CQL), while appearing to solve the problem, actually collapses to behavioral cloning, and that the hard problem was never composition safety but context-awareness.

\section{Problem Setup}

\subsection{Draft Formalization}

A HotS draft consists of 16 sequential actions: 6 bans (removing heroes from the pool) and 10 picks (selecting heroes for teams), alternating between two teams in a fixed order. The state $s_t$ at step $t$ comprises multi-hot vectors for each team's picks and combined bans ($\in \{0,1\}^{90}$ each), map identity (one-hot over 14 maps), skill tier (one-hot over 3 tiers), step number, and step type (ban or pick), for 289 dimensions total. The action space is the set of available heroes ($\sim$84--90 early, decreasing as the draft progresses). Terminal reward is the game outcome (win/loss).

The behavior policy $\pi_\beta$ is defined by 275,285 Storm League replays spanning skill tiers from Bronze through Grandmaster. All models use team-swap augmentation: every sample $(s, y)$ generates a second sample $(s', 1-y)$ by swapping teams and all team-indexed features, doubling the effective training set to 550K.

\subsection{Value Function Variants}

\begin{table}[t]
\centering
\caption{Win probability model variants. All use identical architecture (256$\to$128 MLP, dropout 0.3) and training procedure.}
\label{tab:models}
\begin{tabular}{lccp{3.2cm}}
\toprule
Model & Acc. & Dim & Features \\
\midrule
Naive & 56.5\% & 197 & Multi-hot + map + tier \\
Hero Str. & 57.0\% & 209 & + per-hero WR, team avg \\
Enriched & 57.9\% & 283 & + role counts, pairwise, comp WR, meta \\
\bottomrule
\end{tabular}
\end{table}

Table~\ref{tab:models} shows three win probability models trained on the same data. The enriched model adds 86 features: per-team counts of 9 fine-grained roles (18 dims), pairwise hero counter and synergy statistics (76 dims), and empirical composition win rates (4 dims). The composition WR feature maps each team's 5 heroes to a sorted role tuple and looks up its empirical win rate from aggregate statistics published by a community replay database,\footnote{heroesprofile.com} covering 9.0M games across three skill tiers (121--164 distinct compositions per tier out of 252 possible). Over 99.9\% of training replays match a recorded composition; the 33\% pessimistic default fires for degenerate compositions that humans rarely draft.

\subsection{Evaluation Framework}

We evaluate along two orthogonal axes:

\paragraph{Composition safety.} Healer selection rate, degenerate composition rate (missing healer, missing frontline, no ranged damage, or 3+ heroes in one role). These measure whether the team is structurally functional.

\paragraph{Context-awareness.} \textit{Counter responsiveness}: average normalized pairwise counter delta between our heroes and the opponent's heroes (positive = favorable matchups). \textit{Synergy exploitation}: average normalized within-team synergy delta (positive = heroes that work well together). \textit{Hero diversity}: distinct heroes used, Shannon entropy of the pick distribution, and top-10 hero concentration across all drafts. \textit{GD similarity}: agreement rate with the behavioral cloning model at each pick step.

\paragraph{Cross-model evaluation.} Each draft scored by multiple WP model variants to detect quality trade-offs invisible to single-model metrics.

\subsection{Search Methods}

\textbf{Greedy search:} For each candidate hero, complete the draft via Generic Draft (GD) behavioral cloning opponents trained on 275K replays (289$\to$256$\to$128$\to$90, 5 seed variants for diversity), evaluate the terminal state with $V_\theta$, select the highest-scoring candidate. All $\sim$84 candidates evaluated via batched forward passes.

\textbf{GD behavioral cloning:} Argmax of the GD model at each step. The control for ``what does copying human behavior produce?''

\textbf{CQL:} Argmax of conservatively-learned Q-function. No rollout needed since Q directly estimates action values.

\textbf{MCTS:} AlphaZero-style~\cite{silver2016mastering} with policy and value heads sharing a residual backbone (4.4M parameters). 200 simulations per move, GD opponent sampling during self-play.

\subsection{Sanity Tests}

We evaluate all value functions on 21 hand-crafted scenarios across four categories: absurd compositions (5 tanks, 5 healers, etc.), trap scenarios (high-WR heroes in bad compositions), normal comparisons, and symmetry tests.

\section{The Composition Gap}

\subsection{Cross-Evaluation}

\begin{table}[t]
\centering
\caption{Cross-evaluation matrix. Average WP assigned by column model to row strategy's drafts (480 drafts per strategy). Bold: self-evaluation.}
\label{tab:crosseval}
\begin{tabular}{lccc}
\toprule
& \multicolumn{3}{c}{Evaluated by} \\
\cmidrule(lr){2-4}
Drafter & Naive & Hero Str. & Enriched \\
\midrule
Naive & \textbf{0.670} & 0.595 & 0.586 \\
Hero Str. & 0.619 & \textbf{0.638} & 0.581 \\
Enriched & 0.601 & 0.580 & \textbf{0.644} \\
\bottomrule
\end{tabular}
\end{table}

Table~\ref{tab:crosseval} shows diagonal dominance: each strategy scores highest on its own evaluator, confirming that greedy search optimizes for the specific model's value landscape. The enriched drafter scores lowest on both simpler evaluators (0.601, 0.580), sacrificing 6.9pp of win probability as scored by the naive evaluator to gain compositional quality that only the enriched evaluator perceives.

\subsection{Composition Quality}

\begin{table}[t]
\centering
\caption{Composition quality across 480 drafts per model.\textsuperscript{$\dagger$} All naive-vs-enriched differences significant at $p < 0.0001$ ($\chi^2$ test). Healer difference: 17.5pp, 95\% CI [11.3, 23.7].}
\label{tab:comp}
\begin{tabular}{lccc}
\toprule
Drafter & Healer\% & Ranged\% & Degen.\% \\
\midrule
Naive & 49.0 & 65.2 & 75.8 \\
Hero Str. & 39.0 & 58.1 & 82.1 \\
Enriched & \textbf{66.5} & \textbf{73.3} & \textbf{57.7} \\
\bottomrule
\multicolumn{4}{l}{\footnotesize\textsuperscript{$\dagger$}All models selected frontline heroes $>$95\%.}
\end{tabular}
\end{table}

The hero-strength model performs \textit{worse} than naive on healer selection (39.0\% vs.\ 49.0\%), suggesting individual strength features actively steer the model toward high-WR damage dealers and away from healers, whose individual win rates are moderate. Divergence is high across all draft phases: ban agreement is 9.5\%, early-pick 9.0\%, and late-pick 4.8\%.

\subsection{The Accuracy-Policy Disconnect}

A sweep across 4,100 model configurations ($2^{10}$ subsets of 10 feature groups $\times$ 2 architectures $\times$ 2 learning rates) reveals that the features most important for accuracy differ from those most important for policy quality. Pairwise interaction features dominate accuracy (+4.5\% marginal impact), while the composition WR feature contributes only +0.2\% to accuracy. Yet \texttt{comp\_wr} makes the model assign 17.2\% WP to a 5-tank team (vs.\ 44.7\% from the naive model) and is central to the 18pp reduction in degenerate compositions.

Among the 50 drafts with the largest naive-vs-enriched evaluation gap, 41/50 (82\%) had no healer and 48/50 (96\%) paired individually-strong heroes with degenerate team structure.

\subsection{Independent Baseline}

To verify that OOD overestimation is not specific to our implementation, we reproduce the win probability estimator from Gourdeau's open-source HotS drafter,\footnote{\url{https://github.com/HairyBlob/HotS-Drafter}, the codebase underlying~\cite{gourdeau2021discriminative}.} a shared-weight feedforward network with multiplicative map conditioning using only multi-hot inputs. This baseline (55.3\% accuracy, 13/21 sanity tests) scored 5-tank compositions at 52.8\% WP (rating them as favored), compared to our naive model's 44.6\% and the enriched model's 17.2\%.

\section{Closing the Gap}

\subsection{Residual Gap Analysis}

To provide a controlled baseline, we scale to a 512$\to$256$\to$128 architecture (311K parameters). Even with this additional capacity, the enriched model produces 52\% degenerate compositions. Three factors contribute: search horizon (greedy evaluation cannot reason about future hero availability), training concentration (95\%+ of training teams include a healer, so the model's prior is ``roughly 50\% WP'' for any composition), and feature granularity (composition WR captures role balance but not fine-grained hero-specific interactions).

\subsection{Synthetic Augmentation}

Of the 252 possible 5-role compositions, only 121--164 per tier appear in 9M games. The remaining $\sim$100 compositions per tier are structurally degenerate; no human players draft them. We generate synthetic training records for each unseen composition: heroes sampled by role, paired against real opponent teams, outcomes assigned at a 10\% win rate. 30K synthetic records ($\sim$11\% of training data) are added; the test set remains pure.

We swept 27 configurations: 3 assigned win rates (10\%, 20\%, 30\%) $\times$ 3 volumes (50, 100, 200 per composition per tier) $\times$ 3 scopes (unseen compositions only, sparse low-WR compositions only, or both). Two findings stood out. First, augmenting unseen compositions is the dominant factor; reinforcing sparse-but-real compositions had no effect. Second, a lower assigned win rate (10\%) produces stronger penalties for extreme degeneracy.

\subsection{Augmentation Results}

\begin{table}[t]
\centering
\caption{Effect of synthetic augmentation (200 drafts, 512$\to$256$\to$128 for both rows). Healer difference: 30.5pp ($\chi^2 = 56.6$, $p < 0.0001$). Degenerate difference: 25.5pp ($\chi^2 = 27.3$, $p < 0.0001$).}
\label{tab:synth}
\begin{tabular}{lccccc}
\toprule
Config & Acc.\% & Healer\% & Ranged\% & Degen.\% & Sanity \\
\midrule
No augmentation & 57.2 & 64.0 & 80.5 & 52.0 & 17/21 \\
+ synthetic & 56.4 & \textbf{94.5} & \textbf{88.0} & \textbf{26.5} & \textbf{20/21} \\
\bottomrule
\end{tabular}
\end{table}

The architecture alone (without synthetic data) only reduces degenerate rate from 58\% to 52\% with no healer improvement, confirming synthetic data is the active ingredient. The mechanism is simple: the model learns to trust its own \texttt{comp\_wr} feature because training now includes examples where low composition WR correlates with actual losses.

The residual 26.5\% likely reflects the greedy search horizon (depth-1 cannot plan for future healer availability) and genuine domain ambiguity: no-healer compositions win 43\% of real games.

\section{CQL: Safety Without Understanding}

\subsection{CQL Setup}

We reformulate drafting as offline RL. Each replay provides 16 state-action-reward transitions; with team-swap augmentation this yields $\sim$7.1M training transitions. The Q-network maps state (289 dims) through 512$\to$256$\to$128 layers to 90 Q-values (one per hero), with invalid actions masked to $-\infty$.

We use Monte Carlo returns with $\gamma = 1.0$: since reward is only at the terminal step and episodes are fixed-length, the Q-target for every transition is simply the game outcome. CQL thus becomes supervised Q-regression plus a pessimism penalty on unseen actions:
\begin{equation}
\mathcal{L} = \mathcal{L}_\text{BCE}(Q(s,a), y) + \alpha \left( \mathbb{E}_\mu[\log \sum_a e^{Q(s,a)}] - \mathbb{E}_\mathcal{D}[Q(s,a)] \right)
\end{equation}
where $y$ is the game outcome, $\alpha$ controls conservatism, and the second term penalizes high Q-values for actions not in the data. We sweep $\alpha \in \{0.1, 0.5, 1.0, 2.0, 5.0\}$.

\subsection{The Misleading Result}

\begin{table}[t]
\centering
\caption{CQL composition metrics (200 drafts, naive 289-dim features). CQL appears to subsume feature engineering entirely.}
\label{tab:cql_basic}
\begin{tabular}{lcccc}
\toprule
$\alpha$ & Healer\% & Degen.\% & Front.\% & Ranged\% \\
\midrule
0.1 & 93.5 & 13.0 & 97.0 & 98.0 \\
0.5 & 96.0 & 10.5 & 98.5 & 100.0 \\
1.0 & 93.5 & 9.0 & 98.5 & 99.5 \\
2.0 & 92.0 & 12.5 & 97.5 & 98.0 \\
5.0 & 92.0 & 13.0 & 97.5 & 97.5 \\
\bottomrule
\end{tabular}
\end{table}

Table~\ref{tab:cql_basic} shows CQL with naive multi-hot features achieving 9.0\% degenerate rate ($\alpha = 1.0$), lower than the enriched+augmented model's 26.5\%. Using only hero identity inputs (no compositional features, no synthetic data), CQL appears to completely subsume the feature engineering and augmentation pipeline. This was our initial reading of the results.

\subsection{Rich Evaluation Reveals the Truth}

\begin{table*}[t]
\centering
\caption{Comprehensive evaluation across both axes of draft quality (200 drafts per strategy). Counter and synergy are average normalized pairwise deltas (positive = favorable). Entropy is Shannon entropy of pick distribution. GD Sim is agreement rate with behavioral cloning model. Aug WP is the augmented model's evaluation of each strategy's drafts.}
\label{tab:rich}
\begin{tabular}{lccccccccc}
\toprule
Strategy & Healer\% & Degen.\% & Counter & Synergy & Distinct & Entropy & Top-10\% & GD Sim\% & Aug WP \\
\midrule
GD baseline & 99.5 & 1.0 & $-$0.09 & $-$0.05 & 40 & 4.24 & 75.3 & 83.2 & 0.522 \\
CQL $\alpha\!=\!0.5$ (naive) & 95.0 & 14.5 & $-$0.06 & $-$0.17 & 40 & 4.35 & 73.1 & 56.7 & 0.473 \\
CQL $\alpha\!=\!1.0$ (naive) & 93.5 & 11.0 & $-$0.10 & $-$0.14 & 44 & 4.42 & 70.7 & 61.5 & 0.504 \\
CQL $\alpha\!=\!0.5$ (enriched) & 98.0 & 5.5 & $-$0.10 & $-$0.09 & 38 & 4.24 & 72.3 & -- & -- \\
CQL $\alpha\!=\!2.0$ (enriched) & 98.5 & 3.0 & $-$0.11 & $-$0.04 & 38 & 4.30 & 72.7 & -- & -- \\
\midrule
Enriched WP (greedy) & 77.0 & 42.5 & $+$0.14 & $+$1.03 & 86 & 5.96 & 32.0 & 4.4 & 0.651 \\
Enriched+aug (greedy) & 96.0 & 13.5 & $+$0.06 & $+$0.80 & 90 & 5.97 & 32.8 & 4.5 & -- \\
\bottomrule
\end{tabular}
\end{table*}

Table~\ref{tab:rich} reveals the full picture. We walk through each axis:

\paragraph{Counter responsiveness.} CQL and GD produce negative counter deltas ($-$0.06 to $-$0.11): they do not adapt picks to counter the opponent. The enriched greedy models produce positive deltas (+0.06 to +0.14), indicating genuine matchup-aware selection.

\paragraph{Synergy exploitation.} CQL produces negative synergy ($-$0.04 to $-$0.17). The enriched model produces strongly positive synergy (+1.03 without augmentation, +0.80 with). CQL does not exploit within-team hero interactions.

\paragraph{Hero diversity.} CQL and GD use 38--44 distinct heroes with 70--75\% of picks from the top 10. The enriched models use 86--90 heroes with only 32\% top-10 concentration. CQL's hero pool is collapsed to the same popular subset that humans pick from.

\paragraph{GD similarity.} CQL agrees with the GD behavioral cloning model 57--62\% of the time. The enriched model agrees only 4.4\%. CQL makes decisions highly similar to behavioral cloning; the enriched model makes fundamentally different, context-dependent decisions.

\paragraph{Cross-model WP.} The enriched model scores highest on the augmented WP evaluator (0.651 vs.\ 0.522 for GD and 0.473--0.504 for CQL), confirming it finds higher-quality drafts as perceived by the most compositionally-aware evaluator. CQL scores \textit{lower} than the GD baseline (0.473--0.504 vs.\ 0.522), suggesting CQL degrades draft quality below pure behavioral cloning.

\subsection{CQL with Enriched Features}

Adding the 86 enriched features to CQL's input (375 total dims) improves composition safety (3.0\% degen at $\alpha = 2.0$) but does not produce context-awareness. Counter responsiveness remains negative ($-$0.10 to $-$0.12), synergy remains near-zero ($-$0.04 to $-$0.15), and hero diversity stays at 38 heroes with 72--74\% top-10 concentration. The enriched features help CQL be more precise about what the behavior policy's distribution looks like, but CQL's pessimism still dominates, preventing matchup-dependent decisions.

\subsection{The GD Baseline Revelation}

The most surprising single number in our experiments: the GD behavioral cloning model achieves 99.5\% healer selection and 1.0\% degenerate rate. Composition safety was never the hard problem. Humans in 275K Storm League games virtually always draft balanced teams, so any method that stays near human behavior inherits this property.

CQL's composition safety (93--98\% healer, 3--14\% degen) is \textit{worse} than the GD baseline it approximates, because CQL's Q-value noise causes occasional deviations from the behavior policy in the direction of degenerate compositions. CQL ``solves'' the composition gap not by understanding composition but by approximately reproducing behavioral cloning at the cost of some noise-induced degenerate outputs.

\subsection{Overly Conservative: 100\% Ranged}

CQL $\alpha = 0.5$ achieves 100\% ranged selection, never deviating from the standard template. This is overly conservative: 4-bruiser+healer compositions are genuinely strong in certain matchups. CQL penalizes \textit{all} uncommon actions, not just bad ones, collapsing the hero pool to a safe subset. This is the known trade-off with pessimistic offline RL~\cite{lyu2022mildly} in domains where the behavior policy concentrates on a small subset of the action space.

\section{MCTS with Augmented Value Function}

\subsection{MCTS Training}

We train an AlphaZero-style network (4.4M parameters) via self-play with MCTS. The network shares a residual backbone (289 dims $\to$ 768, 3 residual blocks, compressed to 256) between a perspective-invariant policy head (256 $\to$ 90) and a perspective-aware value head (257 $\to$ 128 $\to$ 64 $\to$ 1, receiving the backbone output plus an our-team indicator). During self-play, the network plays one side with MCTS (200 simulations per move) while the opponent samples from a pool of 5 GD variants.

The value function used for MCTS leaf evaluation varies across runs. We train five configurations with different WP models to measure how value function quality affects MCTS policy quality.

\subsection{MCTS Training Results}

\begin{table}[t]
\centering
\caption{MCTS training runs. WP model used for leaf evaluation during self-play. Win rate and average WP measured vs.\ GD opponents.}
\label{tab:mcts}
\begin{tabular}{llrcc}
\toprule
Run & WP Model & Episodes & WR\% & Avg WP \\
\midrule
A & Base (197d) & 79K & 78.0 & 0.546 \\
C & Enriched (283d) & 167K & 89.0 & 0.598 \\
E & Enriched (300K) & 300K & 89.5 & 0.602 \\
F & + comp\_wr & 300K & 89.5 & 0.598 \\
G & Augmented & 226K$^*$ & 92.5 & 0.626 \\
\bottomrule
\multicolumn{5}{l}{\footnotesize$^*$Run G in progress. Reported at 226K/300K episodes.}
\end{tabular}
\end{table}

Table~\ref{tab:mcts} shows the results. Switching from the base WP model to the enriched model produces an 11.5pp win rate improvement (78.0\% $\to$ 89.5\%). The augmented WP model (Run G) achieves 0.626 average WP at 226K episodes, already higher than any completed run's final WP (0.602), confirming that synthetic augmentation benefits MCTS training through better OOD calibration.

All runs follow the same learning curve: random play at 35--40\% win rate, rapid improvement to 55--60\% by 25K episodes, steady climb to 80--90\% by 100K, and plateau with occasional peaks to 95\%+ after 150K. The learning dynamics are insensitive to the WP model variant; only the asymptote differs.

The comp\_wr feature alone (Run F) does not improve over the enriched model (Run E) in MCTS, likely because the MCTS policy head learns to avoid degenerate compositions through self-play, making the explicit composition penalty redundant. The augmented model's benefit in MCTS comes from better calibration across the full state space, not just from the degenerate composition penalty.

\subsection{Search Depth and Value Function Calibration}

The 26.5\% residual degenerate rate from greedy search is partially a depth-1 limitation. MCTS with 200 simulations provides multi-step lookahead, allowing the agent to plan for future healer availability rather than evaluating one pick at a time. The augmented WP model's higher average WP in MCTS (0.626 vs.\ 0.602) suggests the augmentation helps even with deeper search, by providing more accurate leaf evaluations for the states MCTS explores.

\section{Discussion}

\subsection{Two Orthogonal Axes}

Our results reveal that composition safety and context-awareness require fundamentally different solutions:

\textit{Composition safety} (``does this team have the right roles?'') is trivially solved by copying human behavior. The GD behavioral cloning model achieves 99.5\% healer and 1\% degenerate. CQL provides a noisier approximation of the same mechanism. Neither method requires any understanding of \textit{why} teams need healers; they need only observe that humans always pick them.

\textit{Context-awareness} (``does this team counter the opponent and synergize internally?'') requires value function search with domain-structured features. The enriched model uses 90 distinct heroes (vs.\ 40 for GD/CQL), produces positive counter and synergy deltas, and agrees with GD on only 4.4\% of decisions. It makes fundamentally different picks based on the specific matchup. No behavioral cloning or CQL variant achieves this.

\subsection{Why CQL Collapses}

CQL's pessimism penalizes \textit{all} uncommon state-action pairs, not just bad ones. In draft, the behavior policy is concentrated: humans regularly pick from roughly 40 heroes. CQL's penalty suppresses the other 50 heroes regardless of matchup context. Adding enriched features makes CQL more precise about what ``in-distribution'' looks like (3\% degen vs.\ 11\%) but does not help it distinguish ``uncommon and bad'' from ``uncommon and good for this matchup.''

This is a structural limitation of pessimistic offline RL in domains with concentrated behavior policies, consistent with the findings of Lyu et al.~\cite{lyu2022mildly}. The draft behavior policy's concentration ratio (40 regularly-picked heroes out of 90) is comparable to many real-world domains where the behavior policy reflects conservative institutional norms rather than optimal decision-making.

\subsection{Practical Implications}

For draft recommendation systems, value function search with enriched features plus augmentation produces the best context-aware drafts. CQL and behavioral cloning produce safe but generic recommendations, which may be appropriate for casual players but not for competitive play where matchup-specific adaptation matters.

For value function evaluation, aggregate accuracy is insufficient. Policy-level metrics including counter responsiveness, synergy exploitation, and hero diversity should be standard evaluation criteria for any system that uses learned value functions for search.

For offline RL practitioners, our results caution that CQL's composition metrics can be misleading. Low degenerate rates may indicate behavioral cloning rather than learned compositional reasoning. Action diversity metrics (distinct actions, entropy, behavioral cloning agreement) should be reported alongside reward-based metrics to detect this collapse.

\section{Conclusion}

The composition gap in learned value functions reveals two orthogonal axes of quality in combinatorial selection. Composition safety is trivially achieved by behavioral cloning (99.5\% healer, 1\% degenerate) and approximately achieved by CQL (93--98\% healer, 3--11\% degenerate), but at the cost of zero context-awareness. Context-awareness is achieved only by value function search with domain features (positive counter/synergy, 90 distinct heroes), but requires feature engineering and OOD augmentation to also achieve composition safety.

CQL's apparent success (9\% degenerate with naive features) is misleading: it collapses to behavioral cloning as shown by negative counter/synergy deltas, half the hero diversity, and 62\% GD agreement. Adding enriched features to CQL does not recover context-awareness. MCTS with the augmented value function achieves the highest overall performance (0.626 avg WP, 92.5\% win rate), combining deeper search with a well-calibrated value function.

The solution achieving both axes is domain-structured features (compositional reasoning) + synthetic augmentation (OOD calibration) + value function search (context-aware action selection). This requires no RL algorithm changes: standard supervised learning with targeted training data is sufficient.

\textit{Limitations.} Single game domain, fixed GD opponent model, greedy search for most comparisons. The CQL collapse may be specific to domains with concentrated behavior policies. MCTS Run G in progress at time of writing.

\textit{Future work.} Full rich evaluation metrics for MCTS policies, offline RL methods that preserve action diversity while maintaining conservatism, and testing the safety/context-awareness decomposition in other combinatorial domains (team sports drafts, card game deck building, portfolio selection).

\begin{thebibliography}{20}

\bibitem{kumar2020conservative}
A.~Kumar, A.~Zhou, G.~Tucker, and S.~Levine, ``Conservative Q-Learning for Offline Reinforcement Learning,'' in \textit{Advances in Neural Information Processing Systems}, vol.~33, pp.~1179--1191, 2020.

\bibitem{levine2020offline}
S.~Levine, A.~Kumar, G.~Tucker, and J.~Fu, ``Offline Reinforcement Learning: Tutorial, Review, and Perspectives on Open Problems,'' \textit{arXiv preprint arXiv:2005.01643}, 2020.

\bibitem{kumar2019stabilizing}
A.~Kumar, J.~Fu, M.~Soh, G.~Tucker, and S.~Levine, ``Stabilizing Off-Policy Q-Learning via Bootstrapping Error Reduction,'' in \textit{Advances in Neural Information Processing Systems}, vol.~32, 2019.

\bibitem{chen2018drafting}
Z.~Chen, Y.~Sun, M.~Seif El-Nasr, and T.-H.~D.~Nguyen, ``The Art of Drafting: A Team-Oriented Hero Recommendation System for Multiplayer Online Battle Arena Games,'' in \textit{Proc. ACM Conf. Recommender Systems (RecSys)}, 2018.

\bibitem{gourdeau2021discriminative}
D.~Gourdeau and L.~Archambault, ``Discriminative Neural Network for Hero Selection in Professional Heroes of the Storm and DOTA 2,'' \textit{IEEE Trans. Games}, vol.~13, no.~4, pp.~380--387, 2021.

\bibitem{lee2022draftrec}
H.~Lee, D.~Hwang, H.~Kim, B.~Lee, and J.~Choo, ``DraftRec: Personalized Draft Recommendation for Winning in Multi-Player Online Battle Arena Games,'' in \textit{Proc. ACM Web Conf. (WWW)}, 2022.

\bibitem{summerville2016draft}
A.~Summerville, M.~Cook, and B.~Steenhuisen, ``Draft Prediction for Heroes of the Storm,'' \textit{arXiv preprint arXiv:1601.07803}, 2016.

\bibitem{hanke2017recommending}
L.~Hanke and L.~Chaimowicz, ``A Recommender System for Hero Line-Ups in MOBA Games,'' in \textit{Proc. AAAI Conf. Artificial Intelligence and Interactive Digital Entertainment}, 2017.

\bibitem{pobiedina2013draft}
N.~Pobiedina, J.~Neidhardt, M.~del Carmen Calatrava~Moreno, and H.~Werthner, ``Ranking Factors of Team Success,'' in \textit{Proc. Int. Conf. World Wide Web Companion}, 2013.

\bibitem{lyu2022mildly}
J.~Lyu, X.~Ma, X.~Li, and Z.~Lu, ``Mildly Conservative Q-Learning for Offline Reinforcement Learning,'' in \textit{Advances in Neural Information Processing Systems}, vol.~35, 2022.

\bibitem{silver2016mastering}
D.~Silver \textit{et al.}, ``Mastering the Game of Go with Deep Neural Networks and Tree Search,'' \textit{Nature}, vol.~529, no.~7587, pp.~484--489, 2016.

\bibitem{kalashnikov2018qtopt}
D.~Kalashnikov \textit{et al.}, ``Scalable Deep Reinforcement Learning for Vision-Based Robotic Manipulation,'' in \textit{Conf. Robot Learning}, pp.~651--673, 2018.

\bibitem{segan2026composition}
M.~Segan, E.~Pysher, and J.~Frankel, ``When Similar Accuracy Produces Different Policies: Value Function Features Matter More Than Aggregate Metrics,'' in \textit{Proc. IEEE Conf. Games (CoG)}, 2026.

\end{thebibliography}

\end{document}
